% block: partial-differential-equations-heading page:43 type:section
\subsection{\BT{Partial Differential Equations}{偏微分方程}}\label{sec:partial-differential-equations}
\BT{In this section, the finite difference method and the finite element method are described using some partial differential equations, and we will see that the more accurate solution we need, the larger the linear systems we need to solve.}{本节通过若干偏微分方程介绍有限差分方法和有限元方法, 并将说明: 所需的解越精确, 就需要求解规模越大的线性方程组.}
% block: finite-differences-heading page:43 type:subsection
\subsubsection{\BT{Finite Difference Methods}{有限差分方法}}\label{subsec:finite-differences}
\BT{This subsection describes the finite difference methods using several simple partial differential equations. The notion of tensor products is also explained, which comes in handy to describe the coefficient matrix of the linear systems arising from the finite difference discretization of the partial differential equations. First of all, a mathematical preliminary is given next.}{本小节通过几个简单的偏微分方程介绍有限差分方法. 还将说明张量积的概念, 它便于描述偏微分方程经有限差分离散后所产生线性方程组的系数矩阵. 首先给出一些数学预备知识.}
% block: difference-preliminary-heading page:44 type:subsection
\paragraph{\BT{Preliminary}{预备知识}}\label{subsec:difference-preliminary}
\BT{This subsection describes several formulas for approximating derivatives $f'(x),f''(x)$ of a smooth function $f(x)$ at $x$ using $f(x),f(x-h)$, and $f(x+h)$ with $0<h<1$. Consider how to approximate $f'(x)$. The Taylor series of $f(x)$ is given by}{本小节介绍几个公式, 利用 $f(x),f(x-h)$ 和 $f(x+h)$, 在 $0<h<1$ 时近似光滑函数 $f(x)$ 在 $x$ 处的导数 $f'(x),f''(x)$. 考虑如何近似 $f'(x)$. $f(x)$ 的 Taylor 级数为}
\begin{equation}\label{eq:taylor-difference-expansion}
f(x+h)=f(x)+f'(x)h+\frac{1}{2!}f''(x)h^2+\frac{1}{3!}f'''(x)h^3+\cdots.
\end{equation}
\BT{Then, $f'(x)$ is written as}{于是, $f'(x)$ 可写成}
\[\begin{aligned}
f'(x)&=\frac{f(x+h)-f(x)}{h}-\frac{1}{2!}f''(x)h-\frac{1}{3!}f'''(x)h^2-\cdots\\
&=\frac{f(x+h)-f(x)}{h}+O(h).
\end{aligned}\]
% block: forward-backward-differences page:44 type:ordinary-paragraph
\BT{$O(\cdot)$ is the big O notation. Roughly speaking, $O(h)$ in this case means that $-\frac{1}{2!}f''(x)h-\frac{1}{3!}f'''(x)h^2-\cdots$ can be written as $ch$ for a suitable constant value $c$ as $h\to0$. The approximation to $f'(x)$ by the recurrence}{其中 $O(\cdot)$ 是大 O 记号. 粗略地说, 此处 $O(h)$ 表示当 $h\to0$ 时, 对适当的常数值 $c$, $-\frac{1}{2!}f''(x)h-\frac{1}{3!}f'''(x)h^2-\cdots$ 可写成 $ch$. 用递推关系}
\begin{equation}\label{eq:forward-difference}
f'(x)\approx\frac{f(x+h)-f(x)}{h}
\end{equation}
\BT{is referred to as the forward difference. The accuracy of the forward difference is of order $h$. A similar approximation}{来近似 $f'(x)$, 称为向前差分. 向前差分的精度为 $h$ 阶. 类似的近似}
\begin{equation}\label{eq:backward-difference}
f'(x)\approx\frac{f(x)-f(x-h)}{h}
\end{equation}
\BT{is referred to as the backward difference. The accuracy of the backward difference is also of order $h$. For a more accurate formula, it follows from \eqref{eq:taylor-difference-expansion} that we have}{称为向后差分. 向后差分的精度同样为 $h$ 阶. 为得到更精确的公式, 由 \eqref{eq:taylor-difference-expansion} 有}
\[\frac{f(x+h)-f(x-h)}{2h}=f'(x)+O(h^2),\]
\BT{from which, we have}{由此可得}
\begin{equation}\label{eq:central-first-difference}
f'(x)\approx\frac{f(x+h)-f(x-h)}{2h},
\end{equation}
\BT{and this is referred to as the central difference for the first derivative of $f(x)$, whose accuracy is of order $h^2$.}{这称为 $f(x)$ 的一阶导数的中心差分, 精度为 $h^2$ 阶.}
% block: central-second-difference pages:44,45 type:ordinary-paragraph
\BT{We now consider approximating the second derivative $f''(x)$. From the Taylor series \eqref{eq:taylor-difference-expansion} of $f(x+h)$ and $f(x-h)$, we have}{现在考虑近似二阶导数 $f''(x)$. 由 $f(x+h)$ 和 $f(x-h)$ 的 Taylor 级数 \eqref{eq:taylor-difference-expansion} 有}
\[f(x+h)-2f(x)+f(x-h)=f''(x)h^2+O(h^4).\]
\BT{Thus}{因此}
\[f''(x)=\frac{f(x+h)-2f(x)+f(x-h)}{h^2}+O(h^2).\]
\BT{The following approximation}{下面的近似}
\begin{equation}\label{eq:central-second-difference}
f''(x)\approx\frac{f(x+h)-2f(x)+f(x-h)}{h^2}
\end{equation}
\BT{is referred to as the central difference for the second derivative of $f(x)$, whose accuracy is of order $h^2$.}{称为 $f(x)$ 的二阶导数的中心差分, 精度为 $h^2$ 阶.}
% block: laplace-example-heading page:45 type:subsection
\refstepcounter{applicationexample}
\paragraph{Example~\theapplicationexample\space(\BT{Symmetric Matrix}{对称矩阵})}\label{subsec:laplace-difference}
\begingroup\makeatletter\edef\@currentlabel{\theapplicationexample}\label{ex:laplace-difference}\makeatother\endgroup
\BT{Consider the Laplace equation of the form}{考虑如下形式的 Laplace 方程}
\[\begin{aligned}
u_{xx}+u_{yy}&=0,&(x,y)&\in(0,1)\times(0,1),\\
u(x,y)&=g(x,y),&(x,y)&\in\Gamma,
\end{aligned}\]
\BT{where $\Gamma$ is the boundary of $(0,1)\times(0,1)$. To be specific, the boundary condition is given by $u(0,y)=g(0,y)$, $u(1,y)=g(1,y)$, $u(x,0)=g(x,0)$, and $u(x,1)=g(x,1)$.}{其中 $\Gamma$ 是 $(0,1)\times(0,1)$ 的边界. 具体地, 边界条件为 $u(0,y)=g(0,y)$, $u(1,y)=g(1,y)$, $u(x,0)=g(x,0)$ 和 $u(x,1)=g(x,1)$.}
\BT{We now approximately solve the Laplace equation by using central difference \eqref{eq:central-second-difference}. To this end, an equispaced grid of $N\times N$ points in $(0,1)\times(0,1)$ is used. The case $N=3$ ($h=1/4$) is shown in Fig.\ \ref{fig:laplace-grid}.}{现在用中心差分 \eqref{eq:central-second-difference} 近似求解 Laplace 方程. 为此, 在 $(0,1)\times(0,1)$ 内采用含 $N\times N$ 个点的等距网格. $N=3$ ($h=1/4$) 的情形见 Fig.\ \ref{fig:laplace-grid}.}
% block: laplace-grid-figure page:45 type:figure
\begin{figure}[htbp]\centering
\includegraphics[width=.48\linewidth]{figures/laplace-grid.png}
\caption{\BT{Mesh grid on the region for example\ \ref{ex:laplace-difference} (the Laplace equation)}{Example\ \ref{ex:laplace-difference} (Laplace 方程) 的区域网格}}\label{fig:laplace-grid}
\end{figure}
% block: five-point-discretization page:46 type:ordinary-paragraph
\BT{In the following, we use the mesh in Fig.\ \ref{fig:laplace-grid} and we will see how linear systems are obtained. Central difference \eqref{eq:central-second-difference} for approximating $u_{xx}$ and $u_{yy}$ is given by}{下面使用 Fig.\ \ref{fig:laplace-grid} 中的网格, 说明如何得到线性方程组. 近似 $u_{xx}$ 和 $u_{yy}$ 的中心差分 \eqref{eq:central-second-difference} 为}
\begin{equation}\label{eq:laplace-xx-discretization}
u_{xx}(x,y)\approx\frac{1}{h^2}(u(x+h,y)-2u(x,y)+u(x-h,y)),
\end{equation}
\begin{equation}\label{eq:laplace-yy-discretization}
u_{yy}(x,y)\approx\frac{1}{h^2}(u(x,y+h)-2u(x,y)+u(x,y-h)).
\end{equation}
\BT{It is convenient to introduce an abbreviated notation $u_{i,j}=u(x_i,y_j)$. Then we obtain the following five-point formula:}{引入简写记号 $u_{i,j}=u(x_i,y_j)$ 很方便. 于是得到下面的五点公式:}
\begin{equation}\label{eq:laplace-five-point}
u_{xx}(x_i,y_j)+u_{yy}(x_i,y_j)\approx\frac{1}{h^2}(u_{i+1,j}+u_{i-1,j}-4u_{i,j}+u_{i,j+1}+u_{i,j-1}),
\end{equation}
\BT{where $u(x_i,y_j)=g(x_i,y_j)$ on the boundary. Let $u_1:=u(x_1,y_1)$, $u_2:=u(x_1,y_2)$,\ldots, $u_9:=u(x_3,y_3)$ and $g_{i,j}:=u(x_i,y_j)$ on the boundary. Then from Fig.\ \ref{fig:laplace-grid} and \eqref{eq:laplace-five-point}, unknowns $u_1,u_2,\ldots,u_9$ have the following relations:}{其中在边界上 $u(x_i,y_j)=g(x_i,y_j)$. 令 $u_1:=u(x_1,y_1)$, $u_2:=u(x_1,y_2)$,\ldots, $u_9:=u(x_3,y_3)$, 并在边界上令 $g_{i,j}:=u(x_i,y_j)$. 于是由 Fig.\ \ref{fig:laplace-grid} 和 \eqref{eq:laplace-five-point}, 未知量 $u_1,u_2,\ldots,u_9$ 满足以下关系:}
\[\begin{aligned}
u_2+g_{0,1}-4u_1+u_4+g_{1,0}&=0,\\
u_3+u_1-4u_2+u_5+g_{2,0}&=0,\\
g_{4,1}+u_2-4u_3+u_6+g_{3,0}&=0,\\
u_5+g_{0,2}-4u_4+u_7+u_1&=0,\\
&\vdots\\
g_{4,3}+u_8-4u_9+g_{3,4}+u_6&=0.
\end{aligned}\]
% block: laplace-matrix page:46 type:equation
\BT{The corresponding matrix form is}{相应的矩阵形式为}
\begin{equation}\label{eq:laplace-grid-system}
\underbrace{\left[\begin{array}{rrr|rrr|rrr}
-4&1&0&1&0&0&0&0&0\\1&-4&1&0&1&0&0&0&0\\0&1&-4&0&0&1&0&0&0\\\hline
1&0&0&-4&1&0&1&0&0\\0&1&0&1&-4&1&0&1&0\\0&0&1&0&1&-4&0&0&1\\\hline
0&0&0&1&0&0&-4&1&0\\0&0&0&0&1&0&1&-4&1\\0&0&0&0&0&1&0&1&-4
\end{array}\right]}_A
\underbrace{\begin{bmatrix}u_1\\u_2\\u_3\\u_4\\u_5\\u_6\\u_7\\u_8\\u_9\end{bmatrix}}_{\boldsymbol{x}}
=\underbrace{-\begin{bmatrix}g_{1,0}\\g_{2,0}\\g_{3,0}\\0\\0\\0\\0\\0\\0\end{bmatrix}
-\begin{bmatrix}0\\0\\0\\0\\0\\0\\g_{1,4}\\g_{2,4}\\g_{3,4}\end{bmatrix}
-\begin{bmatrix}g_{0,1}\\0\\0\\g_{0,2}\\0\\0\\g_{0,3}\\0\\0\end{bmatrix}
-\begin{bmatrix}0\\0\\g_{4,1}\\0\\0\\g_{4,2}\\0\\0\\g_{4,3}\end{bmatrix}}_{\boldsymbol{b}}.
\end{equation}
\BT{By solving the linear systems, we obtain the approximate solution $u_1,\ldots,u_9$ on the mesh grid.}{求解这一线性方程组, 即可得到网格上的近似解 $u_1,\ldots,u_9$.}
% block: kronecker-definition pages:46,47 type:ordinary-paragraph
\BT{The matrix of the linear systems \eqref{eq:laplace-grid-system} has a special nonzero structure, which can be rewritten in a simple form using Kronecker product. Let $A(=\{a_{ij}\})$ and $B$ are matrices. Then, the symbol $A\otimes B$ is called the Kronecker product (or Tensor product), defined by}{线性方程组 \eqref{eq:laplace-grid-system} 的矩阵具有特殊的非零结构, 可利用 Kronecker 积将其改写为简单形式. 设 $A(=\{a_{ij}\})$ 和 $B$ 为矩阵. 则符号 $A\otimes B$ 称为 Kronecker 积 (或张量积), 定义为}
\begin{equation}\label{eq:kronecker-product-definition}
A\otimes B:=\begin{bmatrix}a_{11}B&\cdots&a_{1m}B\\\vdots&\ddots&\vdots\\a_{n1}B&\cdots&a_{nm}B\end{bmatrix}.
\end{equation}
\BT{Below is an example for $A,B$ being 2-by-2 matrices:}{下面给出 $A,B$ 为二阶矩阵时的例子:}
\[A\otimes B=\left[\begin{array}{c|c}a_{11}B&a_{12}B\\\hline a_{21}B&a_{22}B\end{array}\right]
=\left[\begin{array}{cc|cc}
a_{11}b_{11}&a_{11}b_{12}&a_{12}b_{11}&a_{12}b_{12}\\
a_{11}b_{21}&a_{11}b_{22}&a_{12}b_{21}&a_{12}b_{22}\\\hline
a_{21}b_{11}&a_{21}b_{12}&a_{22}b_{11}&a_{22}b_{12}\\
a_{21}b_{21}&a_{21}b_{22}&a_{22}b_{21}&a_{22}b_{22}
\end{array}\right].\]
\BT{Now, let}{现在令}
\[M=\begin{bmatrix}-2&1&0\\1&-2&1\\0&1&-2\end{bmatrix},\quad I=\begin{bmatrix}1&0&0\\0&1&0\\0&0&1\end{bmatrix}.\]
\BT{Then, the coefficient matrix $A$ in \eqref{eq:laplace-grid-system} can be written as}{则 \eqref{eq:laplace-grid-system} 中的系数矩阵 $A$ 可写成}
\begin{equation}\label{eq:laplace-kronecker-sum}
A=I\otimes M+M\otimes I.
\end{equation}
\BT{The size of the matrix $A$ is $3^2\times3^2$. In general, $N\times N$ grid yields an $N^2\times N^2$ matrix, which can be very large and symmetric. Krylov subspace methods for symmetric linear systems are described in Sect.\ \ref{sec:hermitian-linear-systems}. Finally, some properties of the Kronecker product are listed below.}{矩阵 $A$ 的规模为 $3^2\times3^2$. 一般而言, $N\times N$ 网格产生一个 $N^2\times N^2$ 矩阵, 它可能非常大, 并且是对称的. 对称线性方程组的 Krylov 子空间方法将在 Section\ \ref{sec:hermitian-linear-systems} 中介绍. 最后, 下面列出 Kronecker 积的一些性质.}
% block: kronecker-properties pages:47,48 type:theorem
\begin{Theorem}\label{thm:kronecker-product-properties}
\BT{The Kronecker product has the following properties:}{Kronecker 积具有以下性质:}
\begin{enumerate}
\item \BT{For $A\in\mathbb{C}^{m\times n},B\in\mathbb{C}^{p\times q},C\in\mathbb{C}^{r\times s}$,}{对 $A\in\mathbb{C}^{m\times n},B\in\mathbb{C}^{p\times q},C\in\mathbb{C}^{r\times s}$,}
\[(A\otimes B)\otimes C=A\otimes(B\otimes C).\]
\item \BT{For $A,B\in\mathbb{C}^{m\times n},C\in\mathbb{C}^{p\times q}$,}{对 $A,B\in\mathbb{C}^{m\times n},C\in\mathbb{C}^{p\times q}$,}
\[(A+B)\otimes C=A\otimes C+B\otimes C.\]
\item \BT{For $A\in\mathbb{C}^{m\times n},B,C\in\mathbb{C}^{p\times q}$,}{对 $A\in\mathbb{C}^{m\times n},B,C\in\mathbb{C}^{p\times q}$,}
\[A\otimes(B+C)=A\otimes B+A\otimes C.\]
\item \BT{For $c\in\mathbb{C},A\in\mathbb{C}^{m\times n},B\in\mathbb{C}^{p\times q}$,}{对 $c\in\mathbb{C},A\in\mathbb{C}^{m\times n},B\in\mathbb{C}^{p\times q}$,}
\[(cA)\otimes B=A\otimes(cB)=c(A\otimes B).\]
\item \BT{For $A\in\mathbb{C}^{m\times n},B\in\mathbb{C}^{p\times q}$,}{对 $A\in\mathbb{C}^{m\times n},B\in\mathbb{C}^{p\times q}$,}
\[(A\otimes B)^{\mathrm H}=A^{\mathrm H}\otimes B^{\mathrm H}.\]
\item \BT{For $A\in\mathbb{C}^{m\times n},B\in\mathbb{C}^{p\times q},C\in\mathbb{C}^{n\times r},D\in\mathbb{C}^{q\times s}$,}{对 $A\in\mathbb{C}^{m\times n},B\in\mathbb{C}^{p\times q},C\in\mathbb{C}^{n\times r},D\in\mathbb{C}^{q\times s}$,}
\[(A\otimes B)(C\otimes D)=(AC)\otimes(BD).\]
\end{enumerate}
\end{Theorem}
% block: kronecker-properties-proof page:48 type:proof
\begin{Proof}
\BT{We only give a proof of (6). The $(i,k)$ block element of $A\otimes B$ is $a_{i,k}B$, and the $(k,j)$ block element of $C\otimes D$ is $c_{k,j}D$. Thus the $(i,j)$ block element of $(A\otimes B)(C\otimes D)$ is given as}{只证明第 (6) 项. $A\otimes B$ 的 $(i,k)$ 块元素为 $a_{i,k}B$, $C\otimes D$ 的 $(k,j)$ 块元素为 $c_{k,j}D$. 因此, $(A\otimes B)(C\otimes D)$ 的 $(i,j)$ 块元素为}
\[\begin{aligned}
((A\otimes B)(C\otimes D))_{i,j}&=\sum_{k=1}^n(a_{i,k}B)(c_{k,j}D)=\left(\sum_{k=1}^na_{i,k}c_{k,j}\right)BD\\
&=(AC)_{i,j}BD=((AC)\otimes(BD))_{i,j},
\end{aligned}\]
\BT{which concludes the proof.}{证明完毕.}
\end{Proof}
% block: block-matrix-meaning page:48 type:ordinary-paragraph
\BT{Note that $(i,j)$ ``block'' element is not a scalar but a matrix. For example, consider the following 2-by-2 block matrices of the form}{注意, $(i,j)$ ``块'' 元素不是标量, 而是矩阵. 例如, 考虑下列形式的二阶分块矩阵}
\[A=\begin{bmatrix}A_{11}&A_{12}\\A_{21}&A_{22}\end{bmatrix},\quad B=\begin{bmatrix}B_{11}&B_{12}\\B_{21}&B_{22}\end{bmatrix}.\]
\BT{Then the $(1,1)$ block element of $A$ is matrix $A_{11}$, and the $(1,2)$ block element of $B$ is matrix $B_{12}$. Furthermore, we can calculate $AB$ as if the submatrices $A_{i,j}$ and $B_{i,j}$ are scalars when the matrix multiplication $A_{i,j}B_{j,k}$ is defined. For example, The $(1,1)$ block element of $AB$ is given by $A_{11}B_{11}+A_{12}B_{21}$. From Theorem\ \ref{thm:kronecker-product-properties}-(6), the following properties hold true.}{于是 $A$ 的 $(1,1)$ 块元素是矩阵 $A_{11}$, $B$ 的 $(1,2)$ 块元素是矩阵 $B_{12}$. 此外, 当矩阵乘法 $A_{i,j}B_{j,k}$ 有定义时, 可以把子矩阵 $A_{i,j}$ 和 $B_{i,j}$ 如同标量一样计算 $AB$. 例如, $AB$ 的 $(1,1)$ 块元素为 $A_{11}B_{11}+A_{12}B_{21}$. 由 Theorem\ \ref{thm:kronecker-product-properties} 的第 (6) 项, 下列性质成立.}
% block: kronecker-inverse page:48 type:corollary
\begin{Corollary}\label{cor:kronecker-inverse}
\BT{Let $A$ and $B$ be invertible matrices. Then}{设 $A$ 和 $B$ 为可逆矩阵. 则}
\[(A\otimes B)^{-1}=A^{-1}\otimes B^{-1}.\]
\end{Corollary}
% block: kronecker-inverse-proof page:48 type:proof
\begin{Proof}
\BT{From Theorem\ \ref{thm:kronecker-product-properties}-(6), we have $(A^{-1}\otimes B^{-1})(A\otimes B)=(A^{-1}A)\otimes(B^{-1}B)=I\otimes I$, which is the Kronecker product of the two identity matrices. It is easy to see that the Kronecker product of the identity matrices is the identity matrix.}{由 Theorem\ \ref{thm:kronecker-product-properties} 的第 (6) 项, 有 $(A^{-1}\otimes B^{-1})(A\otimes B)=(A^{-1}A)\otimes(B^{-1}B)=I\otimes I$, 即两个单位矩阵的 Kronecker 积. 容易看出, 单位矩阵的 Kronecker 积仍是单位矩阵.}
\end{Proof}
% block: kronecker-eigenpairs page:48 type:corollary
\begin{Corollary}\label{cor:kronecker-eigenpairs}
\BT{Let $\boldsymbol{p}_i^{(A)}$ and $\boldsymbol{p}_j^{(B)}$ be eigenvectors of $A$ and $B$, and the corresponding eigenvalues are $\lambda_i^{(A)}$ and $\lambda_j^{(B)}$. Then an eigenvector of $A\otimes B$ is $\boldsymbol{p}_i^{(A)}\otimes\boldsymbol{p}_j^{(B)}$, and the corresponding eigenvalue is $\lambda_i^{(A)}\lambda_j^{(B)}$, i.e.,}{设 $\boldsymbol{p}_i^{(A)}$ 和 $\boldsymbol{p}_j^{(B)}$ 分别为 $A$ 和 $B$ 的特征向量, 对应特征值为 $\lambda_i^{(A)}$ 和 $\lambda_j^{(B)}$. 则 $A\otimes B$ 的一个特征向量为 $\boldsymbol{p}_i^{(A)}\otimes\boldsymbol{p}_j^{(B)}$, 对应特征值为 $\lambda_i^{(A)}\lambda_j^{(B)}$, 即}
\[(A\otimes B)\left(\boldsymbol{p}_i^{(A)}\otimes\boldsymbol{p}_j^{(B)}\right)=\lambda_i^{(A)}\lambda_j^{(B)}\left(\boldsymbol{p}_i^{(A)}\otimes\boldsymbol{p}_j^{(B)}\right).\]
\end{Corollary}
% block: kronecker-eigenpairs-proof pages:48,49 type:proof
\begin{Proof}
\BT{From Theorem\ \ref{thm:kronecker-product-properties}-(4) and (6), we have}{由 Theorem\ \ref{thm:kronecker-product-properties} 的第 (4) 和 (6) 项, 有}
\[(A\otimes B)\left(\boldsymbol{p}_i^{(A)}\otimes\boldsymbol{p}_j^{(B)}\right)=\left(A\boldsymbol{p}_i^{(A)}\right)\otimes\left(B\boldsymbol{p}_j^{(B)}\right)=\lambda_i\lambda_j\left(\boldsymbol{p}_i^{(A)}\otimes\boldsymbol{p}_j^{(B)}\right),\]
\BT{which concludes the proof.}{证明完毕.}
\end{Proof}
% block: kronecker-sum-eigenpairs page:49 type:corollary
\begin{Corollary}\label{cor:kronecker-sum-eigenpairs}
\BT{Using the same notations as in Corollary\ \ref{cor:kronecker-eigenpairs}, it follows that}{沿用 Corollary\ \ref{cor:kronecker-eigenpairs} 中的记号, 有}
\[(A\otimes I+I\otimes B)\left(\boldsymbol{p}_i^{(A)}\otimes\boldsymbol{p}_j^{(B)}\right)=\left(\lambda_i^{(A)}+\lambda_j^{(B)}\right)\left(\boldsymbol{p}_i^{(A)}\otimes\boldsymbol{p}_j^{(B)}\right),\]
\BT{where $I$ is the identity matrix.}{其中 $I$ 为单位矩阵.}
\end{Corollary}
% block: kronecker-sum-invertibility page:49 type:ordinary-paragraph
\BT{The matrix $A\otimes I+I\otimes B$ is referred to as a Kronecker sum (or Tensor sum). It follows from Corollary\ \ref{cor:kronecker-sum-eigenpairs} that eigenvectors and eigenvalues of matrix $A$ in \eqref{eq:laplace-kronecker-sum} can be given by $(\boldsymbol{p}_i\otimes\boldsymbol{p}_j,\lambda_i+\lambda_j)$, where $\boldsymbol{p}_i$ and $\lambda_i$ are the $i$th eigenvector and eigenvalue of $M$. This indicates that matrix $A$ in \eqref{eq:laplace-kronecker-sum} is not invertible if and only if there exist $i,j$ such that $\lambda_i=-\lambda_j$. Thus, the distribution of eigenvalues of much smaller matrices $M$ than $A$ determines whether matrix $A$ is invertible or not.}{矩阵 $A\otimes I+I\otimes B$ 称为 Kronecker 和 (或张量和). 由 Corollary\ \ref{cor:kronecker-sum-eigenpairs} 可知, \eqref{eq:laplace-kronecker-sum} 中矩阵 $A$ 的特征向量和特征值可由 $(\boldsymbol{p}_i\otimes\boldsymbol{p}_j,\lambda_i+\lambda_j)$ 给出, 其中 $\boldsymbol{p}_i$ 和 $\lambda_i$ 为 $M$ 的第 $i$ 个特征向量和特征值. 这表明, \eqref{eq:laplace-kronecker-sum} 中的矩阵 $A$ 不可逆当且仅当存在 $i,j$ 使得 $\lambda_i=-\lambda_j$. 因而, 规模比 $A$ 小得多的矩阵 $M$ 的特征值分布决定了矩阵 $A$ 是否可逆.}
\BT{In general, Kronecker product does not commute, i.e.,}{一般而言, Kronecker 积不满足交换律, 即}
\[A\otimes B\ne B\otimes A,\]
\BT{where $A\in\mathbb{C}^{m\times n}$, $B\in\mathbb{C}^{p\times q}$. However, there are permutation matrices $P_{mp}\in\mathbb{R}^{mp\times mp}$ and $P_{nq}\in\mathbb{R}^{nq\times nq}$ such that}{其中 $A\in\mathbb{C}^{m\times n}$, $B\in\mathbb{C}^{p\times q}$. 但是, 存在置换矩阵 $P_{mp}\in\mathbb{R}^{mp\times mp}$ 和 $P_{nq}\in\mathbb{R}^{nq\times nq}$ 使得}
\[P_{mp}^{\top}(A\otimes B)P_{nq}=B\otimes A.\]
\BT{We describe this fact in detail as given next.}{下面详细说明这一事实.}
% block: kronecker-permutation-theorem pages:49,50 type:theorem
\begin{Theorem}\label{thm:kronecker-permutation}
\BT{Let $P_{mn}\in\mathbb{R}^{mn\times mn}$ be a matrix of the form}{设 $P_{mn}\in\mathbb{R}^{mn\times mn}$ 为如下形式的矩阵}
\begin{equation}\label{eq:kronecker-commutation-matrix}
P_{mn}=\begin{bmatrix}
(E_{11})^{\top}&(E_{12})^{\top}&\cdots&(E_{1n})^{\top}\\
(E_{21})^{\top}&(E_{22})^{\top}&\cdots&(E_{2n})^{\top}\\
\vdots&\vdots&\ddots&\vdots\\
(E_{m1})^{\top}&(E_{m2})^{\top}&\cdots&(E_{mn})^{\top}
\end{bmatrix},
\end{equation}
\BT{where $E_{i,j}\in\mathbb{R}^{m\times n}$ is a matrix whose $(i,j)$ element is one and the other elements are zeros. Then,}{其中 $E_{i,j}\in\mathbb{R}^{m\times n}$ 的 $(i,j)$ 元素为一, 其余元素为零. 则}
\[P_{mn}=P_{nm}^{\top}=P_{nm}^{-1}\]
\BT{and for $A\in\mathbb{C}^{m\times n},B\in\mathbb{C}^{p\times q}$,}{并且对 $A\in\mathbb{C}^{m\times n},B\in\mathbb{C}^{p\times q}$,}
\begin{equation}\label{eq:kronecker-swapping-permutation}
\underbrace{P_{mp}^{\top}}_{mp\times mp}\underbrace{(A\otimes B)}_{mp\times nq}\underbrace{P_{nq}}_{nq\times nq}=B\otimes A.
\end{equation}
\BT{If $A$ and $B$ are square matrices (i.e., $m=n$, $p=q$),}{若 $A$ 和 $B$ 为方阵 (即 $m=n$, $p=q$),}
\[P_{mp}^{\top}(A\otimes B)P_{mp}=B\otimes A.\]
\end{Theorem}
% block: kronecker-permutation-proof page:50 type:proof
\begin{Proof}
\BT{We give a proof of \eqref{eq:kronecker-swapping-permutation}. The $(i,j)$ block element of $P_{m,n}$ in \eqref{eq:kronecker-commutation-matrix} is $E_{ij}^{\top}$, and $E_{ij}$ can be factorized as}{下面证明 \eqref{eq:kronecker-swapping-permutation}. \eqref{eq:kronecker-commutation-matrix} 中 $P_{m,n}$ 的 $(i,j)$ 块元素为 $E_{ij}^{\top}$, 而 $E_{ij}$ 可分解为}
\[E_{ij}=\boldsymbol{e}_i\boldsymbol{e}_j^{\top},\]
\BT{where $\boldsymbol{e}_i$ is the $i$th unit vector, i.e., the $i$th element is one, and the others are zeros. Let $a_{ij}$ and $b_{ij}$ be the $(i,j)$ elements of $A$ and $B$. Then, from \eqref{eq:kronecker-commutation-matrix}, the $(i,j)$ block element of $P_{mp}^{\top}(A\otimes B)P_{nq}$ is calculated as}{其中 $\boldsymbol{e}_i$ 为第 $i$ 个单位向量, 即第 $i$ 个元素为一, 其余元素为零. 设 $a_{ij}$ 和 $b_{ij}$ 分别为 $A$ 和 $B$ 的 $(i,j)$ 元素. 于是, 由 \eqref{eq:kronecker-commutation-matrix}, $P_{mp}^{\top}(A\otimes B)P_{nq}$ 的 $(i,j)$ 块元素计算如下}
\[\begin{aligned}
\left(P_{mp}^{\top}(A\otimes B)P_{nq}\right)_{ij}
&=\sum_{k,\ell}E_{ki}(a_{k\ell}B)E_{\ell j}^{\top}=\sum_{k,\ell}a_{k\ell}\left(E_{ki}BE_{j\ell}\right)\\
&=\sum_{k,\ell}a_{k\ell}\left(\boldsymbol{e}_k\boldsymbol{e}_i^{\top}B\boldsymbol{e}_j\boldsymbol{e}_\ell^{\top}\right)
=\sum_{k,\ell}a_{k\ell}\left(\boldsymbol{e}_kb_{ij}\boldsymbol{e}_\ell^{\top}\right)\\
&=b_{ij}\sum_{k,\ell}a_{k\ell}\left(\boldsymbol{e}_k\boldsymbol{e}_\ell^{\top}\right)
=b_{ij}\left(\sum_{k,\ell}a_{k\ell}E_{k\ell}\right)=b_{ij}A\\
&=(B\otimes A)_{ij},
\end{aligned}\]
\BT{which concludes the proof.}{证明完毕.}
\end{Proof}
% block: helmholtz-example-heading page:50 type:subsection
\refstepcounter{applicationexample}
\paragraph{Example~\theapplicationexample\space(\BT{Complex Symmetric Matrix}{复对称矩阵})}\label{subsec:helmholtz-difference}
\begingroup\makeatletter\edef\@currentlabel{\theapplicationexample}\label{ex:helmholtz-difference}\makeatother\endgroup
\BT{We consider the following boundary value problem:}{考虑下列边值问题:}
\[\begin{aligned}
u_{xx}+u_{yy}+\sigma^2u&=0,&(x,y)&\in(0,\pi)\times(0,\pi),\\
u_x|_{x=0}&=i\sqrt{\sigma^2-\tfrac14}\cos\tfrac y2,&&\text{Neumann Condition (1)},\\
u_x|_{x=\pi}-i\sqrt{\sigma^2-\tfrac14}u|_{x=\pi}&=0,&&\text{Radiation Condition},\\
u_y|_{y=0}&=0,&&\text{Neumann Condition (2)},\\
u|_{y=\pi}&=0,&&\text{Dirichlet Condition},
\end{aligned}\]
\BT{The elliptic partial differential equation is known as the Helmholtz equation, see, e.g., \BookCite{19}. For simplicity, we use a 3-by-3 mesh grid as shown in Figs.\ \ref{fig:helmholtz-grid} and \ref{fig:helmholtz-ghost-grid}. Similar to Example\ \ref{ex:laplace-difference} in the previous section, applying \eqref{eq:laplace-xx-discretization}--\eqref{eq:laplace-five-point} to the Helmholtz equation yields}{这一椭圆型偏微分方程称为 Helmholtz 方程, 例如见 \BookCite{19}. 为简单起见, 使用 Figs.\ \ref{fig:helmholtz-grid} 和 \ref{fig:helmholtz-ghost-grid} 所示的三乘三网格. 与上一小节的 Example\ \ref{ex:laplace-difference} 类似, 将 \eqref{eq:laplace-xx-discretization}--\eqref{eq:laplace-five-point} 应用于 Helmholtz 方程, 得}
\begin{equation}\label{eq:helmholtz-five-point}
-u_{i+1,j}-u_{i-1,j}+(4-h^2\sigma^2)u_{i,j}-u_{i,j+1}-u_{i,j-1}=0.
\end{equation}
\BT{Considering Eq.\ \eqref{eq:helmholtz-five-point} over the mesh points given in Fig.\ \ref{fig:helmholtz-ghost-grid} leads to the following six equations:}{在 Fig.\ \ref{fig:helmholtz-ghost-grid} 给出的网格点上考虑方程 \eqref{eq:helmholtz-five-point}, 得到以下六个方程:}
% block: helmholtz-grid-figure page:51 type:figure
\begin{figure}[htbp]\centering
\includegraphics[width=.48\linewidth]{figures/helmholtz-grid.png}
\caption{\BT{Mesh points on the region}{区域上的网格点}}\label{fig:helmholtz-grid}
\end{figure}
% block: helmholtz-ghost-grid-figure page:51 type:figure
\begin{figure}[htbp]\centering
\includegraphics[width=.48\linewidth]{figures/helmholtz-ghost-grid.png}
\caption{\BT{Mesh points on the region satisfying the boundary conditions}{满足边界条件的区域网格点}}\label{fig:helmholtz-ghost-grid}
\end{figure}
% block: helmholtz-six-equations page:51 type:equation
\begin{equation}\label{eq:helmholtz-six-equations}
\left\{\begin{aligned}
-u_2-g_4-u_4-g_1+(4-h^2\sigma^2)u_1&=0,\\
-u_3-u_1-u_5-g_2+(4-h^2\sigma^2)u_2&=0,\\
-g_5-u_2-u_6-g_3+(4-h^2\sigma^2)u_3&=0,\\
-u_5-g_6-g_8-u_1+(4-h^2\sigma^2)u_4&=0,\\
-u_6-u_4-g_9-u_2+(4-h^2\sigma^2)u_5&=0,\\
-g_7-u_5-g_{10}-u_3+(4-h^2\sigma^2)u_3&=0,
\end{aligned}\right.
\end{equation}
\BT{where $u_1:=u_{1,1}=u(0,0)$, $u_2:=u_{2,1}=u(h,0)$, $u_3:=u_{3,1}=u(2h,0)$, $u_4:=u_{1,2}=u(0,h)$, $u_5:=u_{2,2}=u(h,h)$, $u_6:=u_{3,2}=u(2h,h)$ and $h=\pi/2$. Here $g_k$'s for $k=1,\ldots,10$ in Fig.\ \ref{fig:helmholtz-ghost-grid} can be computed by considering the four boundary conditions as follows: first, from Neumann Condition (2) it follows that}{其中 $u_1:=u_{1,1}=u(0,0)$, $u_2:=u_{2,1}=u(h,0)$, $u_3:=u_{3,1}=u(2h,0)$, $u_4:=u_{1,2}=u(0,h)$, $u_5:=u_{2,2}=u(h,h)$, $u_6:=u_{3,2}=u(2h,h)$, 且 $h=\pi/2$. Fig.\ \ref{fig:helmholtz-ghost-grid} 中 $k=1,\ldots,10$ 对应的 $g_k$ 可通过下述四个边界条件计算: 首先, 由 Neumann Condition (2) 得}
\[\frac{u_4-g_1}{2h}=\frac{u_5-g_2}{2h}=\frac{u_6-g_3}{2h}=0,\]
\BT{and thus}{因此}
\[g_1=u_4,\quad g_2=u_5,\quad g_3=u_6.\]
\BT{Second, from Neumann Condition (1), we have}{其次, 由 Neumann Condition (1) 有}
% block: helmholtz-boundary-elimination page:52 type:ordinary-paragraph
\[u_x(0,0)\approx\frac{u_2-g_4}{2h}=id\cos\left(\frac02\right),\quad
u_x(0,h)\approx\frac{u_5-g_6}{2h}=id\cos\left(\frac12\times\frac h2\right).\]
\BT{or equivalently,}{或者等价地,}
\[g_4=u_2-2idh,\quad g_6=u_5-2idh\cos\left(\frac h4\right),\]
\BT{where $d=\sqrt{\sigma^2-1/4}$. Third, from the Radiation Condition we obtain}{其中 $d=\sqrt{\sigma^2-1/4}$. 第三, 由辐射条件得}
\[\frac{g_5-u_2}{2h}-idu_3=\frac{g_7-u_5}{2h}-idu_6=0\]
\BT{or equivalently,}{或者等价地,}
\[g_5=u_2+2idhu_3,\quad g_7=u_5+2idhu_6.\]
\BT{Finally, it follows from the Dirichlet Condition that $g_8,g_9$ and $g_{10}$ satisfy}{最后, 由 Dirichlet 条件可知, $g_8,g_9$ 和 $g_{10}$ 满足}
\[g_8=g_9=g_{10}=0.\]
% block: helmholtz-unscaled-system page:52 type:equation
\BT{Substituting $g_1,g_2,\ldots,g_{10}$ into the six equations \eqref{eq:helmholtz-six-equations} yields the following linear systems:}{将 $g_1,g_2,\ldots,g_{10}$ 代入六个方程 \eqref{eq:helmholtz-six-equations}, 得到下列线性方程组:}
\begin{equation}\label{eq:helmholtz-unscaled-system}
\begin{bmatrix}
a&-2&0&-2&0&0\\-1&a&-1&0&-2&0\\0&-2&a-idh&0&0&-2\\
-1&0&0&a&-2&0\\0&-1&0&-1&a&-1\\0&0&-1&0&-2&a-idh
\end{bmatrix}
\begin{bmatrix}u_1\\u_2\\u_3\\u_4\\u_5\\u_6\end{bmatrix}
=\begin{bmatrix}f_1\\0\\0\\f_4\\0\\0\end{bmatrix},
\end{equation}
\BT{where $a=\sqrt{\sigma^2-1/4}$, $f_1=-2idh$, and $f_4=-2idh\cos(h/4)$. The coefficient matrix of the linear systems \eqref{eq:helmholtz-unscaled-system} is not symmetric. By scaling each equation, we obtain symmetric linear systems of the form}{其中 $a=\sqrt{\sigma^2-1/4}$, $f_1=-2idh$, $f_4=-2idh\cos(h/4)$. 线性方程组 \eqref{eq:helmholtz-unscaled-system} 的系数矩阵不是对称的. 对每个方程作尺度变换, 得到如下形式的对称线性方程组}
\begin{equation}\label{eq:helmholtz-symmetric-system}
\begin{bmatrix}
a/4&-1/2&0&-1/2&0&0\\-1/2&a/2&-1/2&0&-1&0\\0&-1/2&(a-idh)/4&0&0&-1/2\\
-1/2&0&0&a/2&-1&0\\0&-1&0&-1&a&-1\\0&0&-1/2&0&-1&(a-idh)/2
\end{bmatrix}
\begin{bmatrix}u_1\\u_2\\u_3\\u_4\\u_5\\u_6\end{bmatrix}
=\begin{bmatrix}f_1/4\\0\\0\\f_4/2\\0\\0\end{bmatrix}.
\end{equation}
\BT{The coefficient matrix of the form \eqref{eq:helmholtz-symmetric-system} is non-Hermitian but symmetric, i.e., complex symmetric. Krylov subspace methods for solving complex symmetric linear systems are described in Sect.\ \ref{sec:complex-symmetric-linear-systems}.}{\eqref{eq:helmholtz-symmetric-system} 的系数矩阵不是 Hermitian 的, 但却是对称的, 即复对称的. 求解复对称线性方程组的 Krylov 子空间方法将在 Section\ \ref{sec:complex-symmetric-linear-systems} 中介绍.}
% block: nonsymmetric-example-heading page:53 type:subsection
\refstepcounter{applicationexample}
\paragraph{Example~\theapplicationexample\space(\BT{Nonsymmetric Matrix}{非对称矩阵})}\label{subsec:nonsymmetric-difference}
\begingroup\makeatletter\edef\@currentlabel{\theapplicationexample}\label{ex:nonsymmetric-difference}\makeatother\endgroup
\BT{Consider the following three-dimensional partial differential equation:}{考虑下列三维偏微分方程:}
\[\begin{aligned}
[\boldsymbol{a}\cdot(\nabla*\nabla)+\boldsymbol{b}\cdot\nabla+c]u(x,y,z)&=g(x,y,z),&&\text{in }\Omega,\\
u(x,y,z)&=0,&&\text{on }\partial\Omega.
\end{aligned}\]
\BT{Here, $\Omega$ is the cubic domain, i.e., $(0,1)\times(0,1)\times(0,1)$ and $\partial\Omega$ is the boundary. The symbol ``$*$'' denotes element-wise multiplication. Equivalently, the operator $\boldsymbol{a}\cdot(\nabla*\nabla)+\boldsymbol{b}\cdot\nabla+c$ is rewritten as}{这里 $\Omega$ 为立方体区域, 即 $(0,1)\times(0,1)\times(0,1)$, $\partial\Omega$ 为其边界. 符号 ``$*$'' 表示逐元素乘法. 等价地, 算子 $\boldsymbol{a}\cdot(\nabla*\nabla)+\boldsymbol{b}\cdot\nabla+c$ 可改写为}
\[a_1\frac{\partial^2}{\partial x^2}+a_2\frac{\partial^2}{\partial y^2}+a_3\frac{\partial^2}{\partial z^2}+b_1\frac{\partial}{\partial x}+b_2\frac{\partial}{\partial y}+b_3\frac{\partial}{\partial z}+c\]
\BT{with $\boldsymbol{a}=[a_1,a_2,a_3]^{\top}$, $\boldsymbol{b}=[b_1,b_2,b_3]^{\top}$ and $\boldsymbol{c}=[c_1,c_2,c_3]^{\top}$. Using central difference \eqref{eq:central-second-difference} for the second derivative of $u$ yields}{其中 $\boldsymbol{a}=[a_1,a_2,a_3]^{\top}$, $\boldsymbol{b}=[b_1,b_2,b_3]^{\top}$, $\boldsymbol{c}=[c_1,c_2,c_3]^{\top}$. 对 $u$ 的二阶导数使用中心差分 \eqref{eq:central-second-difference}, 得}
\[\begin{aligned}
a_1u_{xx}(x,y,z)&\approx\frac{a_1}{h^2}(u(x+h,y,z)-2u(x,y,z)+u(x-h,y,z)),\\
a_2u_{yy}(x,y,z)&\approx\frac{a_2}{h^2}(u(x,y+h,z)-2u(x,y,z)+u(x,y-h,z)),\\
a_3u_{zz}(x,y,z)&\approx\frac{a_3}{h^2}(u(x,y,z+h)-2u(x,y,z)+u(x,y,z-h)).
\end{aligned}\]
\BT{Similarly, using central difference \eqref{eq:central-first-difference} for the first derivative of $u$ gives}{类似地, 对 $u$ 的一阶导数使用中心差分 \eqref{eq:central-first-difference}, 得}
\[\begin{aligned}
b_1u_x(x,y,z)&\approx\frac{b_1}{2h}(u(x+h,y,z)-u(x-h,y,z)),\\
b_2u_y(x,y,z)&\approx\frac{b_2}{2h}(u(x,y+h,z)-u(x,y-h,z)),\\
b_3u_z(x,y,z)&\approx\frac{b_3}{2h}(u(x,y,z+h)-u(x,y,z-h)).
\end{aligned}\]
% block: nonsymmetric-grid-operator pages:53,54 type:ordinary-paragraph
\BT{Considering an $(N+1)\times(N+1)\times(N+1)$ grid, together with the central differences above, leads to}{考虑一个 $(N+1)\times(N+1)\times(N+1)$ 网格, 结合上述中心差分, 得到}
\[\begin{aligned}
&[\boldsymbol{a}\cdot(\nabla*\nabla)+\boldsymbol{b}\cdot\nabla+c]u(x,y,z)\\
&\quad\approx\frac{a_1}{h^2}(u_{i+1,j,k}-2u_{i,j,k}+u_{i-1,j,k})+\frac{b_1}{2h}(u_{i+1,j,k}-u_{i-1,j,k})+cu_{i,j,k}\\
&\qquad+\frac{a_2}{h^2}(u_{i,j+1,k}-2u_{i,j,k}+u_{i,j-1,k})+\frac{b_2}{2h}(u_{i,j+1,k}-u_{i,j-1,k})\\
&\qquad+\frac{a_3}{h^2}(u_{i,j,k+1}-2u_{i,j,k}+u_{i,j,k-1})+\frac{b_3}{2h}(u_{i,j,k+1}-u_{i,j,k-1}),
\end{aligned}\]
\BT{where $h=1/(N+1)$. Since $u=0$ on the boundary, following Sect.\ \ref{subsec:laplace-difference}, we obtain}{其中 $h=1/(N+1)$. 由于边界上 $u=0$, 沿用 Section\ \ref{subsec:laplace-difference} 中的方法, 得到}
\begin{equation}\label{eq:three-dimensional-difference-system}
(I_N\otimes I_N\otimes A+I_N\otimes B\otimes I_N+C\otimes I_N\otimes I_N)\boldsymbol{u}=\boldsymbol{g},
\end{equation}
\BT{where $I_N$ is the $N$-by-$N$ identity matrix, each element of $\boldsymbol{u}$ and $\boldsymbol{g}$ corresponds to $u(x_i,y_i,z_i)$ and $g(x_i,y_i,z_i)$, and}{其中 $I_N$ 为 $N$ 行 $N$ 列的单位矩阵, $\boldsymbol{u}$ 和 $\boldsymbol{g}$ 的每个元素分别对应于 $u(x_i,y_i,z_i)$ 和 $g(x_i,y_i,z_i)$, 且}
\[\begin{aligned}
A&:=\frac{a_1}{h^2}M_1+\frac{b_1}{2h}M_2+cI_N,\\
B&:=\frac{a_2}{h^2}M_1+\frac{b_2}{2h}M_2,\\
C&:=\frac{a_3}{h^2}M_1+\frac{b_3}{2h}M_2.
\end{aligned}\]
\BT{Here, $M_1$ and $M_2$ are defined by}{这里 $M_1$ 和 $M_2$ 定义为}
\[M_1:=\begin{bmatrix}-2&1&&&\\1&-2&1&&\\&\ddots&\ddots&\ddots&\\&&1&-2&1\\&&&1&-2\end{bmatrix},\quad
M_2:=\begin{bmatrix}0&1&&&\\-1&0&1&&\\&\ddots&\ddots&\ddots&\\&&-1&0&1\\&&&-1&0\end{bmatrix}.\]
\BT{$M_1$ and $M_2$ are $N$-by-$N$ (tridiagonal Toeplitz) matrices, and blanks in the matrices are zeros.}{其中 $M_1$ 和 $M_2$ 为 $N$ 行 $N$ 列的 (三对角 Toeplitz) 矩阵, 矩阵中的空白位置为零.}
% block: difference-system-size page:54 type:ordinary-paragraph
\BT{The size of linear systems \eqref{eq:three-dimensional-difference-system} is $N^3\times N^3$. If $N=100$, which is an equispaced grid with $h\approx0.01$, then the size of the matrix is a million. If $N=1000$, then the size is now a billion! Therefore, efficient numerical solvers for solving very large (and sparse) linear systems are required for a well approximate solution to the partial differential equation.}{线性方程组 \eqref{eq:three-dimensional-difference-system} 的规模为 $N^3\times N^3$. 若 $N=100$, 即 $h\approx0.01$ 的等距网格, 则矩阵的阶数为一百万. 若 $N=1000$, 阶数就达到十亿! 因此, 要较好地近似求解这一偏微分方程, 就需要高效的数值求解器来求解非常大 (且稀疏) 的线性方程组.}
\BT{The matrix in \eqref{eq:three-dimensional-difference-system} is (real) nonsymmetric if $\boldsymbol{b}\ne0$. Krylov subspace methods for nonsymmetric (non-Hermitian) linear systems are described in Sect.\ \ref{sec:nonhermitian-linear-systems}.}{当 $\boldsymbol{b}\ne0$ 时, \eqref{eq:three-dimensional-difference-system} 中的矩阵是 (实) 非对称的. 非对称 (非 Hermitian) 线性方程组的 Krylov 子空间方法将在 Section\ \ref{sec:nonhermitian-linear-systems} 中介绍.}
% block: finite-elements-heading page:54 type:subsection
\subsubsection{\BT{The Finite Element Method}{有限元方法}}\label{subsec:finite-elements}
\BT{In the previous section, we have seen that a finite difference approximation to (partial) derivatives yields linear recurrence relations, leading to linear systems by using the recurrence relations. The approximate solution corresponds to the values of given grid points, which are the elements of the solution vector for the linear systems. In this section, the Finite Element Method (FEM, hereafter) is introduced, and we will see how linear systems are obtained through a simple example. The advantages of the FEM over the finite difference method are (i) square (or rectangular) grid points on the domain are not necessarily needed, and thus it is suitable for not only square or rectangular but other various domains; (ii) there are many mathematical results for the error analysis based on functional analysis.}{上一小节已说明, 对 (偏) 导数作有限差分近似会得到线性递推关系, 利用这些递推关系即可建立线性方程组. 近似解对应于给定网格点上的值, 这些值就是线性方程组解向量的元素. 本小节介绍有限元方法 (下文简称 FEM), 并通过一个简单例子说明如何得到线性方程组. FEM 相对于有限差分方法的优点是: (i) 不一定需要区域上的正方形 (或矩形) 网格点, 因而它不仅适用于正方形或矩形区域, 也适用于其他各种区域; (ii) 基于泛函分析, 已有许多关于误差分析的数学结果.}
% block: poisson-domain-figure page:55 type:figure
\begin{figure}[htbp]\centering
\includegraphics[width=.48\linewidth]{figures/poisson-domain.png}
\caption{\BT{The domain and the boundary of the Poisson Eqs.\ \eqref{eq:poisson-equation}--\eqref{eq:poisson-neumann}}{Poisson 方程 \eqref{eq:poisson-equation}--\eqref{eq:poisson-neumann} 的区域与边界}}\label{fig:poisson-domain}
\end{figure}
% block: poisson-boundary-problem page:55 type:ordinary-paragraph
\BT{In this section, we consider applying the FEM to the following Poisson equation:}{本小节考虑将 FEM 应用于下列 Poisson 方程:}
\begin{equation}\label{eq:poisson-equation}
-\Delta u(x,y)=f,\quad(x,y)\in\Omega,
\end{equation}
\begin{equation}\label{eq:poisson-dirichlet}
\cdot\ u(x,y)=g,\quad(x,y)\in\Gamma_1,\quad\text{Dirichlet Condition},
\end{equation}
\begin{equation}\label{eq:poisson-neumann}
\cdot\ \frac{\partial}{\partial n}u(x,y)=0,\quad(x,y)\in\Gamma_2,\quad\text{Neumann Condition},
\end{equation}
\BT{where $\Delta:=\partial^2/\partial x^2+\partial^2/\partial y^2$. The domain and the boundary with the boundary conditions are illustrated in Fig.\ \ref{fig:poisson-domain}. In the following, the explanations are based on \BookCite{118}, \BookCite{170}.}{其中 $\Delta:=\partial^2/\partial x^2+\partial^2/\partial y^2$. 区域以及带有边界条件的边界见 Fig.\ \ref{fig:poisson-domain}. 以下说明基于 \BookCite{118}, \BookCite{170}.}
% block: weak-form-heading page:55 type:subsection
\subsubsection{\BT{Weak Form}{弱形式}}\label{subsec:weak-form}
\BT{In this section, an integral equation called the weak form is derived from the Poisson equation. Consider integrating \eqref{eq:poisson-equation} multiplied by an arbitrary function over the domain $\Omega$:}{本小节从 Poisson 方程导出称为弱形式的积分方程. 考虑将 \eqref{eq:poisson-equation} 乘以任意函数后在区域 $\Omega$ 上积分:}
\begin{equation}\label{eq:poisson-weighted-integral}
-\int_\Omega v\Delta u\,\mathrm d\Omega=\int_\Omega vf\,\mathrm d\Omega.
\end{equation}
\BT{If we can find a function $u$ that satisfies \eqref{eq:poisson-weighted-integral} for all $v$, then $u$ is the solution of \eqref{eq:poisson-equation}. The FEM tries to find the solution $u$ that satisfies \eqref{eq:poisson-weighted-integral} instead of solving \eqref{eq:poisson-equation}, which differs from the finite difference methods. Here, function $v$ is referred to as the test function or the weight function.}{若能找到一个对所有 $v$ 都满足 \eqref{eq:poisson-weighted-integral} 的函数 $u$, 则 $u$ 就是 \eqref{eq:poisson-equation} 的解. FEM 试图寻找满足 \eqref{eq:poisson-weighted-integral} 的解 $u$, 而不是直接求解 \eqref{eq:poisson-equation}, 这与有限差分方法不同. 这里函数 $v$ 称为检验函数或权函数.}
% block: gauss-green-weak-form pages:55,56 type:ordinary-paragraph
\BT{Now, in order to simplify the mathematical formulas, let $v$ be an arbitrary smooth function but $v=0$ on $\Gamma_1$. Applying the Gauss--Green theorem to the left-hand side of \eqref{eq:poisson-weighted-integral}, the integration by parts yields}{现在, 为简化数学公式, 令 $v$ 为任意光滑函数, 但在 $\Gamma_1$ 上满足 $v=0$. 对 \eqref{eq:poisson-weighted-integral} 左端应用 Gauss--Green 定理, 分部积分得}
\begin{equation}\label{eq:poisson-integration-by-parts}
\begin{aligned}
\int_\Omega v\Delta u\,\mathrm d\Omega
&=\int_\Gamma v\frac{\partial u}{\partial n}\,\mathrm d\Gamma-\int_\Omega\nabla v\cdot\nabla u\,\mathrm d\Omega\quad\text{(the Gauss--Green theorem)}\\
&=\int_{\Gamma_1}v\frac{\partial u}{\partial n}\,\mathrm d\Gamma_1+\int_{\Gamma_2}v\frac{\partial u}{\partial n}\,\mathrm d\Gamma_2-\int_\Omega\nabla v\cdot\nabla u\,\mathrm d\Omega\\
&=-\int_\Omega\nabla v\cdot\nabla u\,\mathrm d\Omega.
\end{aligned}
\end{equation}
\BT{Here we used the assumption $v=0$ on $\Gamma_1$ and Neumann condition \eqref{eq:poisson-neumann} on $\Gamma_2$, i.e., $\frac{\partial u}{\partial n}=0$ on $\Gamma_2$. Substituting \eqref{eq:poisson-integration-by-parts} into \eqref{eq:poisson-weighted-integral} gives}{这里使用了 $\Gamma_1$ 上 $v=0$ 的假设以及 $\Gamma_2$ 上的 Neumann 条件 \eqref{eq:poisson-neumann}, 即 $\Gamma_2$ 上 $\frac{\partial u}{\partial n}=0$. 将 \eqref{eq:poisson-integration-by-parts} 代入 \eqref{eq:poisson-weighted-integral}, 得}
\[\int_\Omega\nabla v\cdot\nabla u\,\mathrm d\Omega=\int_\Omega vf\,\mathrm d\Omega.\]
\BT{Thus all we have to do is to find $u$ such that}{因此只需寻找 $u$ 使得}
\begin{equation}\label{eq:poisson-weak-form}
\iint_\Omega\left(\frac{\partial v}{\partial x}\frac{\partial u}{\partial x}+\frac{\partial v}{\partial y}\frac{\partial u}{\partial y}\right)\mathrm dx\mathrm dy=\iint_\Omega vf\,\mathrm dx\mathrm dy.
\end{equation}
\BT{The equation does not have the second-order derivative, whereas the original Eq.\ \eqref{eq:poisson-equation} does. In other words, the condition of differentiability of $u$ in \eqref{eq:poisson-weak-form} looks to be weak. Therefore, \eqref{eq:poisson-weak-form} is referred to as the weak form.}{这个方程不含二阶导数, 而原方程 \eqref{eq:poisson-equation} 含有二阶导数. 换言之, \eqref{eq:poisson-weak-form} 对 $u$ 的可微性要求看来较弱. 因此, \eqref{eq:poisson-weak-form} 称为弱形式.}
% block: elementwise-weak-form page:56 type:ordinary-paragraph
\BT{As seen in \eqref{eq:poisson-weak-form}, the weak form is the integral over the whole domain. Now consider that the whole domain is divided into finite elements $e_i$, i.e., $\Omega=\bigcup_{i=1}^n e_i$ and $e_i\cap e_j=\emptyset$ for $i\ne j$. Then we have}{由 \eqref{eq:poisson-weak-form} 可见, 弱形式是在整个区域上积分. 现在考虑将整个区域划分成有限元 $e_i$, 即 $\Omega=\bigcup_{i=1}^n e_i$, 且当 $i\ne j$ 时 $e_i\cap e_j=\emptyset$. 于是有}
\begin{equation}\label{eq:elementwise-poisson-weak-form}
\iint_{e_i}\left(\frac{\partial v}{\partial x}\frac{\partial u}{\partial x}+\frac{\partial v}{\partial y}\frac{\partial u}{\partial y}\right)\mathrm dx\mathrm dy=\iint_{e_i}vf\,\mathrm dx\mathrm dy
\end{equation}
\BT{for $i=1,2,\ldots,n$. An example of the domain decomposition by triangle elements is shown in Fig.\ \ref{fig:triangular-domain-decomposition}. In the next section, we will see how linear systems are obtained from \eqref{eq:elementwise-poisson-weak-form}.}{其中 $i=1,2,\ldots,n$. 用三角形单元划分区域的例子见 Fig.\ \ref{fig:triangular-domain-decomposition}. 下一小节将说明如何从 \eqref{eq:elementwise-poisson-weak-form} 得到线性方程组.}
% block: element-system-derivation-heading page:56 type:subsection
\subsubsection{\BT{Derivation of Linear Systems}{线性方程组的推导}}\label{subsec:element-system-derivation}
\BT{Let finite elements $e_i$ in \eqref{eq:elementwise-poisson-weak-form} be triangular, see Fig.\ \ref{fig:triangular-domain-decomposition}. Each node of the triangular element $e_i$ is given the position and a number, and the node numbers of the triangle are given clockwise, see Fig.\ \ref{fig:triangle-node-order}. As seen in Fig.\ \ref{fig:triangle-node-order}, the position of node $i$ is $(x_i,y_i)$ for $i=1,2,3$.}{令 \eqref{eq:elementwise-poisson-weak-form} 中的有限元 $e_i$ 为三角形, 见 Fig.\ \ref{fig:triangular-domain-decomposition}. 三角形单元 $e_i$ 的每个节点都指定位置和编号, 三角形的节点按顺时针方向编号, 见 Fig.\ \ref{fig:triangle-node-order}. 如 Fig.\ \ref{fig:triangle-node-order} 所示, 节点 $i$ 的位置为 $(x_i,y_i)$, 其中 $i=1,2,3$.}
% block: triangular-domain-decomposition-figure page:57 type:figure
\begin{figure}[htbp]\centering\includegraphics[width=.48\linewidth]{figures/triangular-domain-decomposition.png}
\caption{\BT{The domain decomposition by triangular finite elements}{用三角形有限元划分区域}}\label{fig:triangular-domain-decomposition}\end{figure}
% block: triangle-node-order-figure page:57 type:figure
\begin{figure}[htbp]\centering\includegraphics[width=.4\linewidth]{figures/triangle-node-order.png}
\caption{\BT{The numbering of triangular finite elements}{三角形有限元的编号}}\label{fig:triangle-node-order}\end{figure}
% block: triangular-shape-functions page:57 type:ordinary-paragraph
\BT{Now we introduce the following vector-valued function:}{现在引入下列向量值函数:}
\begin{equation}\label{eq:triangle-shape-functions}
\boldsymbol{h}(x,y)=\begin{pmatrix}h_1(x,y)\\h_2(x,y)\\h_3(x,y)\end{pmatrix}
=\begin{pmatrix}1&1&1\\x_1&x_2&x_3\\y_1&y_2&y_3\end{pmatrix}^{-1}\begin{pmatrix}1\\x\\y\end{pmatrix},
\end{equation}
\BT{and let $\widetilde u$ be an approximate function to the solution function $u$ over the triangular domain such that $\widetilde u$ is a linear function of $x,y$ with $\widetilde u(x_i,y_i)=u_i$ for $i=1,2,3$. Then $\widetilde u$ is given by}{并设 $\widetilde u$ 是三角形区域上解函数 $u$ 的近似函数, $\widetilde u$ 是 $x,y$ 的线性函数, 且对 $i=1,2,3$ 满足 $\widetilde u(x_i,y_i)=u_i$. 则 $\widetilde u$ 为}
\begin{equation}\label{eq:triangle-linear-interpolant}
\widetilde u=\boldsymbol{h}^{\top}\boldsymbol{u}_e=(1\ x\ y)\begin{pmatrix}1&1&1\\x_1&x_2&x_3\\y_1&y_2&y_3\end{pmatrix}^{-\top}\begin{pmatrix}u_1\\u_2\\u_3\end{pmatrix}.
\end{equation}
\BT{Confirming equalities $\widetilde u(x_i,y_i)=u_i$ for $1\le i\le3$ is left as an exercise for the reader.}{等式 $\widetilde u(x_i,y_i)=u_i$ ($1\le i\le3$) 的验证留给读者作为练习.}
% block: element-gradient-integrand pages:57,58 type:ordinary-paragraph
\BT{Now substituting $\widetilde u$ into the left integrand in \eqref{eq:elementwise-poisson-weak-form} and choosing $v$ as the same form of $\widetilde u$, i.e., $v=\boldsymbol{h}^{\top}\boldsymbol{v}_e$ with $\boldsymbol{v}_e:=(v_1,v_2,v_3)^{\top}$ gives}{现在将 $\widetilde u$ 代入 \eqref{eq:elementwise-poisson-weak-form} 左侧的被积函数, 并选择与 $\widetilde u$ 具有相同形式的 $v$, 即 $v=\boldsymbol{h}^{\top}\boldsymbol{v}_e$, 其中 $\boldsymbol{v}_e:=(v_1,v_2,v_3)^{\top}$, 得}
\[\begin{aligned}
\frac{\partial v}{\partial x}\frac{\partial\widetilde u}{\partial x}+\frac{\partial v}{\partial y}\frac{\partial\widetilde u}{\partial y}
&=\frac{\partial}{\partial x}\boldsymbol{v}_e^{\top}\boldsymbol{h}\frac{\partial}{\partial x}\boldsymbol{h}^{\top}\boldsymbol{u}_e+\frac{\partial}{\partial y}\boldsymbol{v}_e^{\top}\boldsymbol{h}\frac{\partial}{\partial y}\boldsymbol{h}^{\top}\boldsymbol{u}_e\\
&=\boldsymbol{v}_e^{\top}\left(\frac{\partial}{\partial x}\boldsymbol{h}\frac{\partial}{\partial x}\boldsymbol{h}^{\top}+\frac{\partial}{\partial y}\boldsymbol{h}\frac{\partial}{\partial y}\boldsymbol{h}^{\top}\right)\boldsymbol{u}_e,
\end{aligned}\]
\BT{where we used the relation $\boldsymbol{h}^{\top}\boldsymbol{v}_e=\boldsymbol{v}_e^{\top}\boldsymbol{h}$. Here, from \eqref{eq:triangle-shape-functions}, it follows that}{其中使用了关系 $\boldsymbol{h}^{\top}\boldsymbol{v}_e=\boldsymbol{v}_e^{\top}\boldsymbol{h}$. 由 \eqref{eq:triangle-shape-functions} 可得}
\begin{equation}\label{eq:triangle-shape-x-derivative}
\boldsymbol{h}_x:=\frac{\partial\boldsymbol{h}}{\partial x}
=\begin{pmatrix}1&1&1\\x_1&x_2&x_3\\y_1&y_2&y_3\end{pmatrix}^{-1}\begin{pmatrix}0\\1\\0\end{pmatrix}
=\frac1D\begin{pmatrix}y_2-y_3\\y_3-y_1\\y_1-y_2\end{pmatrix}:=\begin{pmatrix}b_1\\b_2\\b_3\end{pmatrix},
\end{equation}
\begin{equation}\label{eq:triangle-shape-y-derivative}
\boldsymbol{h}_y:=\frac{\partial\boldsymbol{h}}{\partial y}
=\begin{pmatrix}1&1&1\\x_1&x_2&x_3\\y_1&y_2&y_3\end{pmatrix}^{-1}\begin{pmatrix}0\\0\\1\end{pmatrix}
=\frac1D\begin{pmatrix}x_3-x_2\\x_1-x_3\\x_2-x_1\end{pmatrix}:=\begin{pmatrix}c_1\\c_2\\c_3\end{pmatrix},
\end{equation}
\BT{where $D$ is defined by}{其中 $D$ 定义为}
\[D:=\det\begin{pmatrix}1&1&1\\x_1&x_2&x_3\\y_1&y_2&y_3\end{pmatrix}.\]
% block: element-stiffness-matrix page:58 type:ordinary-paragraph
\BT{$D$ is twice the area of the triangle. Since $\boldsymbol{h}_x$ and $\boldsymbol{h}_y$ do not depend on $x,y$, the left side of \eqref{eq:elementwise-poisson-weak-form} is calculated as follows:}{其中 $D$ 是三角形面积的两倍. 由于 $\boldsymbol{h}_x$ 和 $\boldsymbol{h}_y$ 不依赖于 $x,y$, \eqref{eq:elementwise-poisson-weak-form} 的左端计算如下:}
\[\begin{aligned}
\iint_{e_i}\left(\frac{\partial v}{\partial x}\frac{\partial\widetilde u}{\partial x}+\frac{\partial v}{\partial y}\frac{\partial\widetilde u}{\partial y}\right)\mathrm dx\mathrm dy
&=\boldsymbol{v}_e^{\top}(\boldsymbol{h}_x\boldsymbol{h}_x^{\top}+\boldsymbol{h}_y\boldsymbol{h}_y^{\top})\boldsymbol{u}_e\iint_{e_i}\mathrm dx\mathrm dy\\
&=\boldsymbol{v}_e^{\top}\frac{(\boldsymbol{h}_x\boldsymbol{h}_x^{\top}+\boldsymbol{h}_y\boldsymbol{h}_y^{\top})D}{2}\boldsymbol{u}_e\\
&=\boldsymbol{v}_e^{\top}A_e\boldsymbol{u}_e.
\end{aligned}\]
\BT{Here, $\iint_{e_i}\mathrm dx\mathrm dy$ is the area of triangle $e_i$, i.e., $D/2$. Matrix $A_e$ is a $3\times3$ symmetric matrix whose $(i,j)$ element, $a_{ij}^{(e)}$, is given by}{这里 $\iint_{e_i}\mathrm dx\mathrm dy$ 为三角形 $e_i$ 的面积, 即 $D/2$. 矩阵 $A_e$ 是一个 $3\times3$ 对称矩阵, 其 $(i,j)$ 元素 $a_{ij}^{(e)}$ 由下式给出}
\begin{equation}\label{eq:element-stiffness-entry}
a_{ij}^{(e)}=\frac D2(b_ib_j+c_ic_j),
\end{equation}
\BT{from \eqref{eq:triangle-shape-x-derivative} and \eqref{eq:triangle-shape-y-derivative}. The right-hand side of \eqref{eq:elementwise-poisson-weak-form} is written as}{这是由 \eqref{eq:triangle-shape-x-derivative} 和 \eqref{eq:triangle-shape-y-derivative} 得到的. \eqref{eq:elementwise-poisson-weak-form} 的右端可写成}
\begin{equation}\label{eq:element-load-vector}
\boldsymbol{f}_e:=\begin{pmatrix}f_1^{(e)}\\f_2^{(e)}\\f_3^{(e)}\end{pmatrix}
=\begin{pmatrix}\iint_{e_i}f_0h_1\,\mathrm dx\mathrm dy\\\iint_{e_i}f_0h_2\,\mathrm dx\mathrm dy\\\iint_{e_i}f_0h_3\,\mathrm dx\mathrm dy\end{pmatrix}
=\frac{f_0D}{6}\begin{pmatrix}1\\1\\1\end{pmatrix},
\end{equation}
\BT{where we used $f(x,y)=f_0$ (constant function) for simplicity. Then, we have}{其中为简单起见, 使用了 $f(x,y)=f_0$ (常函数). 于是有}
\[\iint_{e_i}vf\,\mathrm dx\mathrm dy=\boldsymbol{v}_e^{\top}\boldsymbol{f}_e.\]
% block: element-assembly-principle page:59 type:ordinary-paragraph
\BT{On calculating the right-hand side of \eqref{eq:element-load-vector}, the integrand $h_i$ ($1\le i\le3$) is a triangular pyramid with the height being 1 and the area of the base being $D/2$. Thus the volume of the triangle pyramid is $(D/2)\cdot1\cdot1/3=D/6$. From this, the left-hand side and the right-hand side of \eqref{eq:elementwise-poisson-weak-form} are approximated by}{计算 \eqref{eq:element-load-vector} 的右端时, 被积函数 $h_i$ ($1\le i\le3$) 为高为 1, 底面积为 $D/2$ 的三棱锥. 因此三棱锥的体积为 $(D/2)\cdot1\cdot1/3=D/6$. 由此, \eqref{eq:elementwise-poisson-weak-form} 的左端和右端分别近似为}
\begin{equation}\label{eq:element-weak-matrix-pair}
\boldsymbol{v}_e^{\top}A_e\boldsymbol{u}_e,\quad\boldsymbol{v}_e^{\top}\boldsymbol{f}_e.
\end{equation}
\BT{The above process is the derivation of \eqref{eq:element-weak-matrix-pair} regarding a small area $e_i$. Considering all the linear systems for all $e_i$'s together with boundary conditions leads to large linear systems. By solving the large linear systems, the approximate solution $u_i$ on each node can be obtained. In the next subsection, we will see how large linear systems can be obtained.}{上述过程是对小区域 $e_i$ 推导 \eqref{eq:element-weak-matrix-pair}. 将所有 $e_i$ 的线性方程组与边界条件一同考虑, 就得到大型线性方程组. 求解该大型线性方程组, 即可得到各节点上的近似解 $u_i$. 下一小节将说明如何得到大型线性方程组.}
% block: finite-element-example-heading page:59 type:subsection
\subsubsection{Example}\label{subsec:finite-element-example}
\BT{As an example of the previous subsection, we consider Poisson Eq.\ \eqref{eq:poisson-equation} on the $(0,1)\times(0,1)$ square domain $\Omega$ in Fig.\ \ref{fig:square-poisson-domain}, and we will see how linear systems are obtained.}{作为上一小节的例子, 考虑 Fig.\ \ref{fig:square-poisson-domain} 中 $(0,1)\times(0,1)$ 正方形区域 $\Omega$ 上的 Poisson 方程 \eqref{eq:poisson-equation}, 并说明如何得到线性方程组.}
\BT{We first decompose the square domain in Fig.\ \ref{fig:square-poisson-domain} into eight triangles in Fig.\ \ref{fig:eight-triangle-mesh}. Next, calculate \eqref{eq:element-weak-matrix-pair} for $A_e$ and $\boldsymbol{f}_e$ in each triangle element in Fig.\ \ref{fig:eight-triangle-mesh}, and finally, combine all the information, leading to linear systems.}{首先把 Fig.\ \ref{fig:square-poisson-domain} 中的正方形区域划分成 Fig.\ \ref{fig:eight-triangle-mesh} 中的八个三角形. 接着, 对 Fig.\ \ref{fig:eight-triangle-mesh} 中每个三角形单元的 $A_e$ 和 $\boldsymbol{f}_e$ 计算 \eqref{eq:element-weak-matrix-pair}, 最后将所有信息组合起来, 得到线性方程组.}
\BT{In what follows, we calculate \eqref{eq:element-weak-matrix-pair} for $A_e$ and $\boldsymbol{f}_e$. We see from Fig.\ \ref{fig:eight-triangle-mesh} that there are two kinds of triangles, i.e., type I for $e_1^{\mathrm I},e_3^{\mathrm I},e_5^{\mathrm I},e_7^{\mathrm I}$ and type II for $e_2^{\mathrm {II}},e_4^{\mathrm {II}},e_6^{\mathrm {II}},e_8^{\mathrm {II}}$, leading to two types of linear systems corresponding to \eqref{eq:laplace-five-point}. The details of the two triangles are shown in Fig.\ \ref{fig:two-triangle-types}.}{下面对 $A_e$ 和 $\boldsymbol{f}_e$ 计算 \eqref{eq:element-weak-matrix-pair}. 由 Fig.\ \ref{fig:eight-triangle-mesh} 可见, 有两类三角形: $e_1^{\mathrm I},e_3^{\mathrm I},e_5^{\mathrm I},e_7^{\mathrm I}$ 为 I 型, $e_2^{\mathrm {II}},e_4^{\mathrm {II}},e_6^{\mathrm {II}},e_8^{\mathrm {II}}$ 为 II 型, 从而产生与 \eqref{eq:laplace-five-point} 对应的两类线性方程组. 两类三角形的详情见 Fig.\ \ref{fig:two-triangle-types}.}
% block: square-poisson-domain-figure page:59 type:figure
\begin{figure}[htbp]\centering\includegraphics[width=.48\linewidth]{figures/square-poisson-domain.png}
\caption{\BT{The $(0,1)\times(0,1)$ square domain}{正方形区域 $(0,1)\times(0,1)$}}\label{fig:square-poisson-domain}\end{figure}
% block: eight-triangle-mesh-figure page:60 type:figure
\begin{figure}[htbp]\centering\includegraphics[width=.48\linewidth]{figures/eight-triangle-mesh.png}
\caption{\BT{Triangle mesh}{三角形网格}}\label{fig:eight-triangle-mesh}\end{figure}
% block: two-triangle-types-figure page:60 type:figure
\begin{figure}[htbp]\centering\includegraphics[width=.55\linewidth]{figures/two-triangle-types.png}
\caption{\BT{Two types of triangles}{两类三角形}}\label{fig:two-triangle-types}\end{figure}
% block: type-one-element pages:59,60 type:ordinary-paragraph
\BT{For the type I triangle, it follows from \eqref{eq:triangle-shape-x-derivative}, \eqref{eq:triangle-shape-y-derivative}, and recalling $D$ being twice the area of the triangle (i.e., $D=h^2$), we have}{对 I 型三角形, 由 \eqref{eq:triangle-shape-x-derivative}, \eqref{eq:triangle-shape-y-derivative}, 并回顾 $D$ 为三角形面积的两倍 (即 $D=h^2$), 有}
\[\begin{pmatrix}b_1\\b_2\\b_3\end{pmatrix}=\frac1{h^2}\begin{pmatrix}-h\\h\\0\end{pmatrix},\quad
\begin{pmatrix}c_1\\c_2\\c_3\end{pmatrix}=\frac1{h^2}\begin{pmatrix}-h\\0\\h\end{pmatrix}.\]
\BT{Then, from \eqref{eq:element-stiffness-entry} and \eqref{eq:element-load-vector}, we obtain}{于是由 \eqref{eq:element-stiffness-entry} 和 \eqref{eq:element-load-vector} 得}
\begin{equation}\label{eq:type-one-element-matrix}
A_{e^{\mathrm I}}=\frac12\begin{pmatrix}2&-1&-1\\-1&1&0\\-1&0&1\end{pmatrix},\quad
\boldsymbol{f}_{e^{\mathrm I}}=\frac{f_0h^2}{6}\begin{pmatrix}1\\1\\1\end{pmatrix},
\end{equation}
\BT{which correspond to \eqref{eq:element-weak-matrix-pair}. Similarly, for the type II triangle, we have}{它们对应于 \eqref{eq:element-weak-matrix-pair}. 类似地, 对 II 型三角形有}
\[\begin{pmatrix}b_1\\b_2\\b_3\end{pmatrix}=\frac1{h^2}\begin{pmatrix}0\\h\\-h\end{pmatrix},\quad
\begin{pmatrix}c_1\\c_2\\c_3\end{pmatrix}=\frac1{h^2}\begin{pmatrix}-h\\h\\0\end{pmatrix},\]
\BT{and thus}{因而}
\begin{equation}\label{eq:type-two-element-matrix}
A_{e^{\mathrm {II}}}=\frac12\begin{pmatrix}1&-1&0\\-1&2&-1\\0&-1&1\end{pmatrix},\quad
\boldsymbol{f}_{e^{\mathrm {II}}}=\frac{f_0h^2}{6}\begin{pmatrix}1\\1\\1\end{pmatrix}.
\end{equation}
% block: local-global-node-table page:61 type:table
\begin{table}[htbp]\centering
\caption{\BT{Relation between local node numbers and global node numbers.}{局部节点编号与全局节点编号的关系.}}\label{tab:local-global-node-numbers}
\begin{tabular}{c|cccccccc}\hline
\BT{Element}{单元}&$e_1^{\mathrm I}$&$e_2^{\mathrm {II}}$&$e_3^{\mathrm I}$&$e_4^{\mathrm {II}}$&$e_5^{\mathrm I}$&$e_6^{\mathrm {II}}$&$e_7^{\mathrm I}$&$e_8^{\mathrm {II}}$\\\hline
\shortstack{\BT{Node number of}{三角形单元}\\\BT{triangle element}{节点编号}}&\multicolumn{8}{c}{\BT{Whole node numbers}{整体节点编号}}\\\hline
1&1&2&2&3&4&5&5&6\\2&2&5&3&6&5&8&6&9\\3&4&4&5&5&7&7&8&8\\\hline
\end{tabular}\end{table}
% block: first-element-assembly page:61 type:ordinary-paragraph
\BT{In the following, we explain how to combine all the information to obtain linear systems. In order to combine them, it is required to relate global node numbers 1--9 in Fig.\ \ref{fig:eight-triangle-mesh} with local node numbers 1, 2, and 3 in Fig.\ \ref{fig:two-triangle-types}, and the resulting table is shown in Table\ \ref{tab:local-global-node-numbers}. Finally, linear systems are derived as follows. We first consider triangle element $e_1^{\mathrm I}$. From \eqref{eq:type-one-element-matrix} and Table\ \ref{tab:local-global-node-numbers}, we have}{下面说明如何组合所有信息来得到线性方程组. 为此, 需要将 Fig.\ \ref{fig:eight-triangle-mesh} 中的全局节点编号 1--9 与 Fig.\ \ref{fig:two-triangle-types} 中的局部节点编号 1, 2, 3 对应起来, 所得结果见 Table\ \ref{tab:local-global-node-numbers}. 最后按如下方式导出线性方程组. 首先考虑三角形单元 $e_1^{\mathrm I}$. 由 \eqref{eq:type-one-element-matrix} 和 Table\ \ref{tab:local-global-node-numbers} 有}
\[A_{e_1^{\mathrm I}}=\frac12\begin{array}{c@{}c}
\begin{pmatrix}2&-1&-1\\-1&1&0\\-1&0&1\end{pmatrix}&\begin{matrix}1\\2\\4\end{matrix}\\
\begin{matrix}1&2&4\end{matrix}&\end{array},\quad
\boldsymbol{f}_{e^{\mathrm I}}=\frac{f_0h^2}{6}\begin{array}{c@{}c}\begin{pmatrix}1\\1\\1\end{pmatrix}&\begin{matrix}1\\2\\4\end{matrix}\end{array}\]
\BT{and by using the information we construct the following matrix and right-hand side:}{利用这些信息, 构造如下矩阵和右端项:}
\[\frac12\begin{array}{c@{}c}
\begin{pmatrix}
\mathbf2&\mathbf{-1}&0&\mathbf{-1}&0&0&0&0&0\\
\mathbf{-1}&\mathbf1&0&\mathbf0&0&0&0&0&0\\
0&0&0&0&0&0&0&0&0\\
\mathbf{-1}&\mathbf0&0&\mathbf1&0&0&0&0&0\\
0&0&0&0&0&0&0&0&0\\0&0&0&0&0&0&0&0&0\\0&0&0&0&0&0&0&0&0\\0&0&0&0&0&0&0&0&0\\0&0&0&0&0&0&0&0&0
\end{pmatrix}&\begin{matrix}\mathbf1\\\mathbf2\\3\\\mathbf4\\5\\6\\7\\8\\9\end{matrix}\\
\begin{matrix}\mathbf1&\mathbf2&3&\mathbf4&5&6&7&8&9\end{matrix}&
\end{array},\quad\frac{f_0h^2}{6}\begin{array}{c@{}c}
\begin{pmatrix}\mathbf1\\\mathbf1\\0\\\mathbf1\\0\\0\\0\\0\\0\end{pmatrix}&\begin{matrix}\mathbf1\\\mathbf2\\3\\\mathbf4\\5\\6\\7\\8\\9\end{matrix}
\end{array}\]
\BT{for the linear systems. We next consider triangle element $e_2^{\mathrm {II}}$. From \eqref{eq:type-two-element-matrix} and Table\ \ref{tab:local-global-node-numbers}, we have}{作为线性方程组的矩阵和右端项. 接着考虑三角形单元 $e_2^{\mathrm {II}}$. 由 \eqref{eq:type-two-element-matrix} 和 Table\ \ref{tab:local-global-node-numbers} 有}
\[A_{e_2^{\mathrm {II}}}=\frac12\begin{array}{c@{}c}
\begin{pmatrix}1&-1&0\\-1&2&-1\\0&-1&1\end{pmatrix}&\begin{matrix}2\\5\\4\end{matrix}\\
\begin{matrix}2&5&4\end{matrix}&\end{array},\quad
\boldsymbol{f}_{e^{\mathrm {II}}}=\frac{f_0h^2}{6}\begin{array}{c@{}c}\begin{pmatrix}1\\1\\1\end{pmatrix}&\begin{matrix}2\\5\\4\end{matrix}\end{array}\]
% block: second-element-assembly page:62 type:ordinary-paragraph
\BT{and adding the information to the aforementioned matrix and the right-hand side yields}{将这些信息加到前述矩阵和右端项上, 得到}
\[\frac12\begin{array}{c@{}c}
\begin{pmatrix}
2&-1&0&-1&0&0&0&0&0\\
-1&\mathbf2&0&\mathbf0&\mathbf{-1}&0&0&0&0\\
0&0&0&0&0&0&0&0&0\\
-1&\mathbf0&0&\mathbf2&\mathbf{-1}&0&0&0&0\\
0&\mathbf{-1}&0&\mathbf{-1}&\mathbf2&0&0&0&0\\
0&0&0&0&0&0&0&0&0\\0&0&0&0&0&0&0&0&0\\0&0&0&0&0&0&0&0&0\\0&0&0&0&0&0&0&0&0
\end{pmatrix}&\begin{matrix}1\\\mathbf2\\3\\\mathbf4\\\mathbf5\\6\\7\\8\\9\end{matrix}\\
\begin{matrix}1&\mathbf2&3&\mathbf4&\mathbf5&6&7&8&9\end{matrix}&
\end{array},\quad\frac{f_0h^2}{6}\begin{array}{c@{}c}
\begin{pmatrix}1\\\mathbf2\\0\\\mathbf2\\\mathbf1\\0\\0\\0\\0\end{pmatrix}&\begin{matrix}1\\\mathbf2\\3\\\mathbf4\\\mathbf5\\6\\7\\8\\9\end{matrix}
\end{array}\]
\BT{and similarly for the remaining elements $e_3^{\mathrm I},e_4^{\mathrm {II}},\ldots,e_8^{\mathrm {II}}$, we finally have}{对其余单元 $e_3^{\mathrm I},e_4^{\mathrm {II}},\ldots,e_8^{\mathrm {II}}$ 作类似处理, 最终得到}
\[\frac12\begin{pmatrix}
2&-1&0&-1&0&0&0&0&0\\-1&4&-1&0&-2&0&0&0&0\\0&-1&2&0&0&-1&0&0&0\\
-1&0&0&4&-2&0&-1&0&0\\0&-2&0&-2&8&-2&0&-2&0\\0&0&-1&0&-2&4&0&0&-1\\
0&0&0&-1&0&0&2&-1&0\\0&0&0&0&-2&0&-1&4&-1\\0&0&0&0&0&-1&0&-1&2
\end{pmatrix}\begin{pmatrix}u_1\\u_2\\u_3\\u_4\\u_5\\u_6\\u_7\\u_8\\u_9\end{pmatrix}
=\frac{f_0h^2}{6}\begin{pmatrix}1\\3\\2\\3\\6\\3\\2\\3\\1\end{pmatrix}.\]
% block: finite-element-dirichlet-elimination page:62 type:ordinary-paragraph
\BT{The linear systems do not include the information of Dirichlet condition ($u(x,0)=0$ for $x\in[0,1]$). It follows from Figs.\ \ref{fig:square-poisson-domain} and \ref{fig:eight-triangle-mesh} that the following condition is required for satisfying the Dirichlet condition:}{这一线性方程组尚未包含 Dirichlet 条件 (对 $x\in[0,1]$, 有 $u(x,0)=0$) 的信息. 由 Figs.\ \ref{fig:square-poisson-domain} 和 \ref{fig:eight-triangle-mesh} 可知, 满足 Dirichlet 条件要求}
\[u_1=u_2=u_3=0.\]
\BT{Thus, removing the first three rows from the linear systems and substituting $u_1=u_2=u_3=0$ into the resulting linear systems yields:}{因此, 从线性方程组中删去前三行, 并将 $u_1=u_2=u_3=0$ 代入所得线性方程组, 得到:}
\begin{equation}\label{eq:finite-element-square-system}
\frac12\begin{pmatrix}
4&-2&0&-1&0&0\\-2&8&-2&0&-2&0\\0&-2&4&0&0&-1\\
-1&0&0&2&-1&0\\0&-2&0&-1&4&-1\\0&0&-1&0&-1&2
\end{pmatrix}\begin{pmatrix}u_4\\u_5\\u_6\\u_7\\u_8\\u_9\end{pmatrix}
=\frac{f_0h^2}{6}\begin{pmatrix}3\\6\\3\\2\\3\\1\end{pmatrix}.
\end{equation}
\BT{By solving the linear systems we have $u_4,u_5,\ldots,u_9$.}{求解线性方程组, 即可得到 $u_4,u_5,\ldots,u_9$.}
% block: finite-element-exact-comparison page:63 type:ordinary-paragraph
\BT{Finally, we compare the approximate solution with the exact solution. The exact solution of the Poisson equation is $u(x,y)=f_0y-\frac12f_0y^2$. When $f_0=1$, the values of $u_4,u_5,u_6$ for the exact solution are 0.375 and $u_7,u_8,u_9$ are 0.5. Solving linear systems \eqref{eq:finite-element-square-system} with $f_0=1$ and $h=\frac12$ yields $u_4\approx0.381,u_5\approx0.375,u_6\approx0.369$ and $u_7\approx0.524,u_8\approx0.500,u_9\approx0.476$.}{最后比较近似解与精确解. Poisson 方程的精确解为 $u(x,y)=f_0y-\frac12f_0y^2$. 当 $f_0=1$ 时, 精确解的 $u_4,u_5,u_6$ 的值为 0.375, $u_7,u_8,u_9$ 的值为 0.5. 取 $f_0=1$, $h=\frac12$ 求解线性方程组 \eqref{eq:finite-element-square-system}, 得 $u_4\approx0.381,u_5\approx0.375,u_6\approx0.369$ 和 $u_7\approx0.524,u_8\approx0.500,u_9\approx0.476$.}
\BT{If a more accurate approximate solution is required, the number of nodes should be increased. From the above derivation, we see that the size of the linear systems is almost the same as the number of the whole nodes.}{若需要更精确的近似解, 应增加节点数. 由上述推导可见, 线性方程组的规模与整体节点数几乎相同.}
